\section{The Fibonacci normal form and two sphericality tests}\label{sec:fib}
\begin{lemma}[Fibonacci normal form]\label{fib:lem:normal-form}
A rank-four fusion ring containing $f$ with $f^2=1+f$ is self-dual
and has a basis $1,f,X,Y$ with
\begin{equation}\label{fib:eq:fib-auto-rules}
\begin{gathered}
 fX=Y,\qquad fY=X+Y,\qquad X^2=1+aX+bY,\\
 XY=f+bX+(a+b)Y,\qquad
 Y^2=1+f+(a+b)X+(a+2b)Y,
\end{gathered}
\end{equation}
where $a,b\ge0$ are integers. Its positive character is
$(1,\phi,x,\phi x)$, where $\phi=(1+\sqrt5)/2$ and
$x^2=1+(a+b\phi)x$. In a normalized signed sphericalization quotient,
the same formulas hold with integers $\alpha,\beta$ satisfying
$|\alpha|\le a$, $|\beta|\le b$, $\alpha\equiv a$ and
$\beta\equiv b\pmod2$.
\end{lemma}
\begin{proof}
On the two simples outside the Fibonacci subring, left and right
multiplication by $f$ give commuting symmetric nonnegative integral
matrices satisfying $M^2=I+M$. Their trace is one and determinant is
$-1$, so each is a reordering of
$\left(\begin{smallmatrix}0&1\\1&1\end{smallmatrix}\right)$.
The two reorderings do not commute; hence the left and right actions
coincide. Duality commutes with this action and cannot exchange the
two outside simples. All four simples are therefore self-dual.
Frobenius reciprocity and associativity give the displayed products.
For a signed quotient, first change the sign of $f$ so that
$f^2=1+f$. A sign change of an outside basis element normalizes its
nonzero off-diagonal action coefficient to one. The same calculation
then gives the signed formulas; domination and parity give the stated
restrictions on $\alpha,\beta$.
\end{proof}

\begin{proposition}\label{fib:prop:small-spherical}
Every complex categorification of either Fibonacci normal form
$(a,b)=(1,0)$ or $(a,b)=(1,1)$ admits a spherical structure.
\end{proposition}
\begin{proof}
The first table has multiplicity one, so
Corollary~\ref{core:cor:multiplicity-one-spherical} applies.
For $(a,b)=(1,1)$, the signed parameters $\alpha,\beta$ belong to
$\{1,-1\}$. Reversing the signs of both outside basis elements changes
both parameters' signs. Equal signs therefore recover the ordinary
table up to gauge and give a spherical section.

For opposite signs it suffices to take
$t=\alpha+\beta\phi=1-\phi$. A character evaluating $f$ at $\phi$
and $X$ at a root of $x^2-tx-1$ has character field
\[
 L=\Q(\sqrt5,\sqrt{6-\phi}),\qquad
 N_{\Q(\sqrt5)/\Q}(6-\phi)=29.
\]
The radicand is not a square in $\Q(\sqrt5)$, since its norm
is not a rational square. If $L/\Q$ were normal, its two
conjugate quadratic extensions of $\Q(\sqrt5)$ would agree.
Their radicands would have square ratio, forcing $29$ to be a square
in $\Q(\sqrt5)$. This is impossible. Thus $L$ is not normal,
contrary to cyclotomicity of signed character values from
Lemma~\ref{core:lem:family-induction}. Only the ordinary gauge remains.
\end{proof}

